\documentclass{beamer}

\usepackage{graphicx,hyperref,fes,url}
\usepackage[english]{babel}
\usepackage[latin1]{inputenc}
%\usepackage{times}
\usepackage[T1]{fontenc}
\usepackage{multirow}
\usepackage{capt-of}
\usepackage{graphicx}
\usepackage{array}
\usepackage{tikz}
\usepackage{hyperref}


% The title of the presentation:
%  - first a short version which is visible at the bottom of each slide;
%  - second the full title shown on the title slide;
\title[FES style for Beamer]{
  UAlberta FES style for Beamer \LaTeX}

% Optional: a subtitle to be dispalyed on the title slide

% The author(s) of the presentation:
%  - again first a short version to be displayed at the bottom;
%  - next the full list of authors, which may include contact information;

\author[Author Short]{
  Pim Vullers MSc \and Martijn Sprengers \\\medskip
  {\small \url{p.vullers@cs.ru.nl} \and \url{m.sprengers@student.ru.nl}} \\
  {\small \url{http://www.cs.ru.nl/~pim/}}}

% The institute:
%  - to start the name of the university as displayed on the top of each slide
%    this can be adjusted such that you can also create a Dutch version
%  - next the institute information as displayed on the title slide
\institute[University of Alberta]{
  Future Energy Systems Symposium\\
  University of Alberta}

% Add a date and possibly the name of the event to the slides
%  - again first a short version to be shown at the bottom of each slide
%  - second the full date and event name for the title slide

\renewcommand{\(}{\begin{columns}}
\renewcommand{\)}{\end{columns}}
\newcommand{\<}[1]{\begin{column}{#1}}
\renewcommand{\>}{\end{column}}
%%%%%%%%%%%%%%%%%%%%%%%%%%%%%%%%%%%%%%%%%%%%%%%%%%





\newcommand{\degC}{$^o$C$\,$}
\newcommand{\co}{$\text{CO}_{2}\,$}
\newcommand{\tx}{$G_{2\times\text{CO}_2}$}
\newcommand{\Real}{\mathbb R}



% The main document



\begin{document}




\title[FES Symposium Presentation]{Carbon prices in Alberta's Power Market}
\subtitle{And a few other things we're working on...}

%\author{Branko Bo\v skovi\'{c} \and Andrew Leach \and Charles F.\ Mason}
%\institute{University of Alberta  \and University of Alberta  \and University of Wyoming}

%\author[shortname]{author1 \inst{1} \and author2 \inst{2}}
%\institute[shortinst]{\inst{1} affiliation for author1 \and %
 %                     \inst{2} affiliation for author2}

\author[Theme 6, Grids and Storage]{Theme 6: Grids and Storage\
\\\small{Presentation by Andrew Leach}\\\small{Storage content provided by Marc Secanell\\Smart Grid content provided by Ryan Li}}%
 % \texorpdfstring{%
 %   \begin{columns}
%      \column{.1\linewidth}
%      \column{.4\linewidth}
%      \centering
%      Andrew Leach \\University of Alberta
%      \column{.4\linewidth}
%      \centering
%      Blake Shaffer \\University of Calgary
%      \column{.1\linewidth}
%      \end{columns}
% }
% {Leach and Shaffer (2019)}
%}


\date[May, 2019]{May, 2019}

\frame
  \titlepage


\section{Introduction}
\begin{frame}[t]{Grids and Storage Theme Summary}
\begin{block}{Goal} Develop advanced smart grid and energy storage technologies and study the economics and policy to enable the high penetration of renewable energy into the grid.
\end{block}
\begin{block}{Active projects: Smart Grids}
\begin{itemize}\setlength\itemsep{.01em}
\item Future Smart Grid Structures
\item Operational Decision Support for Smart Grid
\item Energy Management Strategies for Distributed Energy Resources
\item Surface Science Hub for Clean Technology
\end{itemize}
\end{block}
\end{frame}

\begin{frame}[t]{Grids and Storage Theme Summary}
\begin{block}{Goal} Develop advanced smart grid and energy storage technologies and study the economics and policy to enable the high penetration of renewable energy into the grid.
\end{block}
\begin{block}{Active projects:Storage}
\begin{itemize}\setlength\itemsep{.01em}
\item Distributed Energy Storage
\item Utility-Scale Energy Storage
\item Next Generation Energy Storage Using Unconventional Materials
\item Research on Interfacial Control of Solid State Lithium Batteries
\end{itemize}

\end{block}
\end{frame}

\begin{frame}[t]{Grids and Storage Theme Summary}
\begin{block}{Goal} Develop advanced smart grid and energy storage technologies and study the economics and policy to enable the high penetration of renewable energy into the grid.
\end{block}
\begin{block}{Active projects: Policies \& Markets}
\begin{itemize}\setlength\itemsep{.01em}
\item Economic policy and the future electricity grid
\item Carbon pricing and other policies impacts on the market for storage
\item The complementary roles of renewables and storage on the grid
\item Computing Platform for Managing Transactive Microgrids
\end{itemize}
\end{block}
\end{frame}



\section{Policies and Markets}
\begin{frame}[t]{Carbon Pricing, Electricity Markets and Output-Based Allocations}
\vbox to \textheight{
  \begin{itemize}
    \item Impacts of carbon pricing has become a central policy question in Canada and globally
    \item Electricity markets, at least in certain jurisdictions are ideal environments in which to study responses to carbon prices
    \item Rich, hourly electricity market data allows us to observe firm-level decisions with respect to carbon price pass through
    \item Fabra and Reguant (2014) shows how firms respond to emissions credit prices and that lump-sum credit allocations were not distortionary in the short run
    \item Brown et al. (2018) looks at Alberta policy changes which alter both carbon prices and output-based credit allocations
    \end{itemize}
  \vfill}
\end{frame}


%\begin{frame}{Oil sands reserves\hspace{3cm} \small Source: AER (2013)}
%\pgfputat{\pgfxy(0,.25)}{\pgfimage[width=4.25in]{AB_reserves.pdf}}
%\end{frame}

\begin{frame}{Alberta's Power Market is an Ideal Laboratory}
\vbox to \textheight{
  \begin{itemize}
    \item Alberta's power market is relatively small, isolated, and features a mix of generation technologies
    \item Our peak internal load for 2018 was 11,697 MW, while our lowest observed internal load was 	7,819 MW
    \item We have relatively minimal intertie capacity.  Alberta is connected to WECC (1325 MW export and 1500 MW import path-rating) via Montana and BC) and to Saskatchewan (150 MW path-rating) but significant hourly variability in capacity which helps us with identification
    \item Supply mix (2018) was 44.8\% coal, 21.1\% net-to-grid from combined heat and power plants, 18.5\% natural gas, 2.9\% hydro, 6.4\% wind, 	1.0\% other sources including solar and 5.2\% imports.
  \end{itemize}
  \vfill}
\end{frame}

\begin{frame}{Alberta's Policy Changes as Treatments}
\vbox to \textheight{
    \begin{itemize}
    \item From 2007 to 2015, Alberta's \textit{Specified Gas Emitters Regulation} imposed a carbon price of \$15 per tonne with output-based allocations at a rate of 88\% of a facility's historic emissions intensity if the facility's annual emissions were more than 100,000 tonnes per year
    \item In 2015, these parameters were changed to \$20 per tonne with allocations equal to 85\% of historic emissions intensities in 2016 and \$30 per tonne with allocations of 80\% in 2017.
    \item 2007-2017 system also had generous emissions credits for combined heat and power systems
    \item In 2018, the system changed again, to one with output-based allocations fixed for all generators (0.37t/MWh) including for net generation from combined heat and power.
    \end{itemize}
\vfill}
\end{frame}


\begin{frame}{Data Richness}
\vbox to \textheight{
  \begin{block}{We have data on:}
    \begin{itemize}
    \item Hourly merit order (power offers) by unit with data on who held offer control for that unit at that time
    \item Hourly forecast and actual loads and prices
    \item Hourly weather for Edmonton, Calgary, and Fort McMurray
    \item Hourly imports, exports, and intertie capabilities
    \item Hourly wind forecasts including 3 and 7 day advanced forecasts.
    \item Plant-level emissions intensities and historic allocation rates under \textit{SGER}
    \end{itemize}
  \end{block}
\vfill}
\end{frame}


\begin{frame}{Carbon Pricing Policy and Power Generation Sources}
   \tikz [remember picture,overlay]
    \node[yshift=-1.15cm,xshift=0cm] at (current page.center)
       {\includegraphics[width=.75\paperwidth]{../../alberta_power/gen_ghg_price.png}}; \vspace{1cm}
   \vfill
\end{frame}


\begin{frame}{Carbon Intensity of Alberta Power}
   \tikz [remember picture,overlay]
    \node[yshift=-1.75cm,xshift=0cm] at (current page.center)
       {\includegraphics[width=.9\paperwidth]{../../alberta_power/monthly_ghg_mwh.png}}; \vspace{1cm}
   \vfill
\end{frame}

\begin{frame}{Alberta Power Prices and Carbon Policy}
   \tikz [remember picture,overlay]
    \node[yshift=-1.75cm,xshift=0cm] at (current page.center)
       {\includegraphics[width=.85\paperwidth]{../../alberta_power/forwards_long.png}}; \vspace{1cm}
   \vfill
\end{frame}



\section{Storage Technology}
\begin{frame}[t]{Secanell group}
\footnotesize{\begin{itemize} \setlength\itemsep{.01em}
  \item Design, procurement and installation of low and high temperature electrolysis fabrication and testing facilities
  \item Fabrication of single-cell low temperature electrolyzer prototypes using iridium and iridium oxide catalysts
  \item Demonstrated, for the first time, the use of inkjet printing as a feasible technology to fabricate low temperature electrolyzer cells  (https://doi.org/10.1149/2.1101807jes)
   \end{itemize}}
\end{frame}


\begin{frame}[t]{Secanell Group}
   \tikz [remember picture,overlay]
    \node[yshift=-.5cm,xshift=0cm] at (current page.center)
      {\includegraphics[height=2.7in]{../../secanell_1}}; \vspace{1cm}
   \vfill
\end{frame}

\begin{frame}[t]{Secanell Group}
\footnotesize{\begin{itemize} \setlength\itemsep{.01em}
  \item Design, procurement and installation of low and high temperature electrolysis fabrication and testing facilities
  \item Fabrication of single-cell low temperature electrolyzer prototypes using iridium and iridium oxide catalysts
  \item Demonstrated, for the first time, the use of inkjet printing as a feasible technology to fabricate low temperature electrolyzer cells  (https://doi.org/10.1149/2.1101807jes)
   \item Development of a novel nano-pompon-like iridium catalyst superstructure for low temperature fuel cell electrolyzers (https://doi.org/10.1016/j.jcat.2019.01.018)
   \end{itemize}}
   \vfill
\end{frame}

\begin{frame}[t]{Secanell Group}
   \tikz [remember picture,overlay]
    \node[yshift=-1.05cm,xshift=0cm] at (current page.center)
      {\includegraphics[height=2.7in]{../../secanell_2}}; \vspace{1cm}
   \vfill
\end{frame}

\begin{frame}[t]{Secanell Group}
\footnotesize{\begin{itemize} \setlength\itemsep{.01em}
  \item Design, procurement and installation of low and high temperature electrolysis fabrication and testing facilities
  \item Fabrication of single-cell low temperature electrolyzer prototypes using iridium and iridium oxide catalysts
  \item Demonstrated, for the first time, the use of inkjet printing as a feasible technology to fabricate low temperature electrolyzer cells  (https://doi.org/10.1149/2.1101807jes)
   \item Development of a novel nano-pompon-like iridium catalyst superstructure for low temperature fuel cell electrolyzers (https://doi.org/10.1016/j.jcat.2019.01.018)
   \item Development of a novel technique for characterization of iridium catalyst degradation (https://doi.org/10.1016/j.jcat.2019.01.018)
 \item Development of novel micro-structures for high temperature solid oxide electrolyzers and fuel cells (https://doi.org/10.1016/j.electacta.2018.02.055)
  \end{itemize}}
   \vfill
\end{frame}



\begin{frame}[t]{Secanell Group}
   \tikz [remember picture,overlay]
    \node[yshift=-1.05cm,xshift=0cm] at (current page.center)
      {\includegraphics[width=4.5in]{../../secanell_3}}; \vspace{1cm}
   \vfill
\end{frame}

\begin{frame}[t]{Secanell Group}
\footnotesize{\begin{itemize} \setlength\itemsep{.01em}
  \item Design, procurement and installation of low and high temperature electrolysis fabrication and testing facilities
  \item Fabrication of single-cell low temperature electrolyzer prototypes using iridium and iridium oxide catalysts
  \item Demonstrated, for the first time, the use of inkjet printing as a feasible technology to fabricate low temperature electrolyzer cells  (https://doi.org/10.1149/2.1101807jes)
   \item Development of a novel nano-pompon-like iridium catalyst superstructure for low temperature fuel cell electrolyzers (https://doi.org/10.1016/j.jcat.2019.01.018)
   \item Development of a novel technique for characterization of iridium catalyst degradation (https://doi.org/10.1016/j.jcat.2019.01.018)
 \item Development of novel micro-structures for high temperature solid oxide electrolyzers and fuel cells (https://doi.org/10.1016/j.electacta.2018.02.055)
  \end{itemize}
HQP include 5 PhD students and 1 MSc, 5 faculty (1 in MECE, 3 in CHEME and 1 in Chem). We are also collaborating with Ionomr, a firm in BC (https://ionomr.com/).}
   \vfill
\end{frame}


\section{Smart Grids}
\begin{frame}{Smart Grids}
Most renewable energy based power generation, energy storage systems and modern electronic loads are based on DC technologies. This necessitates new grid structure such as hybrid AC/DC and DC grids
\begin{block}{Main objectives}
\begin{itemize}\setlength\itemsep{.01em}
\item Development of hybrid AC-DC grids technologies
\item ower electronics technologies, and power management technology
\end{itemize}
HQP include PI: Ryan Li, co-PI: Khajehoddin, Kish, Mohamed, Salmon, Liang. 18 total HQP since 2017
\end{block}
\end{frame}
\begin{frame}
\begin{block}{Main progress so far}
\begin{itemize}\setlength\itemsep{.01em}
\item New DC/AC and DC/DC converter technologies for future grids
\item Power management, power quality management
\item Over 60 publications and 8 new IP
\end{itemize}
\end{block}

\begin{block}{Collaborators}
\begin{itemize}\setlength\itemsep{.01em}
\item NSERC CREATE initiative with RWTH Aachen
\item Collaboration with Tsinghua Univ
\item Alberta utility companies
\item Manufacturers: Rockwell Automation, Landmark,
\item Alberta CTD projects on ``Smart Gas Field Microgrids" and ``Next-Gen Modular Green Building Energy Systems With Optional Energy Storage"
\end{itemize}
\end{block}

\end{frame}

\begin{frame}{Smart Grid Lab}
   \tikz [remember picture,overlay]
    \node[yshift=-1.05cm,xshift=0cm] at (current page.center)
      {\includegraphics[width=4.5in]{../../li_lab}}; \vspace{1cm}
   \vfill
\end{frame}



\begin{frame}{Contact info}
\begin{center}
Andrew Leach\bigskip \\
School of Business, University of Alberta\bigskip \\
Email: \href{mailto:aleach@ualberta.ca}{andrew.leach@ualberta.ca}\bigskip \\
Twitter: \href{http://twitter.com/andrew_leach}{\url{@andrew_leach}}
\end{center}
\vfill
\end{frame}


\end{document}